\section{Methodology}

\subsection{Hypothesis: Preference Entropy Collapse Beyond Performance-Justified Levels}

We hypothesize that extended RLHF training (beyond initial performance saturation) induces preference entropy collapse faster than task performance improves. Specifically:

\textbf{Core Statement:} Under RLHF training on preference datasets, if models undergo extended alignment optimization (10K-20K training steps), then preference distribution entropy collapses faster than task performance improves (inflection point where dH/dP accelerates), because users habituate to model output style and converge on preferences even for subjective tasks where diversity is legitimate.

\textbf{Causal Mechanism (4 steps):}
\begin{enumerate}
\item \textbf{Base models} $\rightarrow$ high-variance outputs $\rightarrow$ diverse preferences $\rightarrow$ \textbf{high entropy} (baseline state)
\item \textbf{Early RLHF} (1K-5K steps) $\rightarrow$ quality improvement $\rightarrow$ entropy decreases \textbf{proportionally} to performance (justified reduction)
\item \textbf{Extended RLHF} (10K-20K steps) $\rightarrow$ performance saturates $\rightarrow$ entropy reduction \textbf{accelerates} beyond performance gains (inflection point)
\item \textbf{Task stratification} $\rightarrow$ subjective tasks show \textbf{greater entropy collapse} than objective tasks (monoculture evidence)
\end{enumerate}

\textbf{Key Assumptions:}
\begin{itemize}
\item \textbf{A1:} Preference entropy proxies critical evaluation capacity (higher H = diverse judgments)
\item \textbf{A5:} Existing RLHF datasets contain raw preference distributions (not just aggregated win rates)
\end{itemize}

This paper focuses on validating \textbf{Assumption A5} via sub-hypothesis H-E1. The full causal mechanism (steps 2-4) requires this prerequisite to succeed.

\subsection{Sub-Hypothesis H-E1: Preference Entropy Existence}

\textbf{Statement:} Under RLHF preference datasets (Anthropic-HH), if we have access to raw pairwise comparison data, then we can compute Shannon entropy H = -$\Sigma$ $p_i$ log($p_i$) for preference distributions, because the dataset publishes response frequency counts rather than only aggregated win rates.

\textbf{Verification Protocol:}
\begin{enumerate}
\item Download Anthropic-HH dataset (160K+ pairwise comparisons, public repository)
\item Sample n=100 prompts with $\geq$5 comparisons each (seed=1 for reproducibility)
\item Group examples by prompt, aggregate preference distributions
\item Compute Shannon entropy H for each prompt's preference distribution
\item Validate entropy range [0, ln(2)] for binary comparisons (0 = perfect consensus, ln(2) $\approx$ 0.693 = maximum uncertainty)
\item Report success rate (\% of prompts with computable entropy) and variance (std of entropy values)
\end{enumerate}

\textbf{Success Criteria (MUST\_WORK gate):}
\begin{itemize}
\item \textbf{Primary:} Entropy computable for $\geq$95\% of sampled prompts
\item \textbf{Secondary:} Entropy variance > 0 (not all constant values)
\end{itemize}

\textbf{Rationale:} If entropy cannot be computed from existing datasets (A5 violated), the entire hypothesis mechanism becomes untestable without new data collection. H-E1 validates the foundational measurement assumption before investing in full causal chain experiments.

\subsection{Entropy Measurement}

\textbf{Shannon Entropy Definition:}

\begin{equation}
H = -\sum_{i=1}^{n} p_i \log(p_i)
\end{equation}

where:
\begin{itemize}
\item $p_i$ = proportion of preferences for option $i$
\item $n$ = number of response options per prompt
\item Natural log (base $e$) for units in nats (vs log$_2$ for bits)
\end{itemize}

\textbf{For Binary Pairwise Comparisons:}
\begin{itemize}
\item $n = 2$ (chosen vs rejected)
\item Theoretical range: $[0, \ln(2)] \approx [0, 0.693]$ nats
\item Maximum entropy (uniform): $H = \ln(2)$ when $p_{\text{chosen}} = p_{\text{rejected}} = 0.5$
\item Minimum entropy (consensus): $H = 0$ when all prefer one option
\end{itemize}

\textbf{Expected Behavior:}
\begin{itemize}
\item High diversity prompts (e.g., opinion questions): $H > 0.4$ nats
\item Low diversity prompts (e.g., factual QA with clear answer): $H < 0.3$ nats
\item Variance across prompts: std$(H) > 0.1$ nats indicates meaningful spread
\end{itemize}

\subsection{Dataset: Anthropic-HH}

\textbf{Source:} Anthropic Helpful \& Harmless (HH) dataset \cite{bai2022constitutional}
\begin{itemize}
\item \textbf{Repository:} https://github.com/anthropics/hh-rlhf
\item \textbf{Size:} 160,800 training examples (helpful-base: 122K, harmless-base: 38K)
\item \textbf{Format:} Each example contains:
  \begin{itemize}
  \item \texttt{chosen}: full conversation text (Human + Assistant turns) for preferred response
  \item \texttt{rejected}: full conversation text for less-preferred response
  \end{itemize}
\item \textbf{Collection:} Single-annotator pairwise comparisons (one labeler per example)
\end{itemize}

\textbf{Sampling Strategy:}
\begin{itemize}
\item Sample n=100 prompts from training split (seed=1)
\item Filter: require $\geq$5 examples per prompt group (ensure sufficient data for entropy estimation)
\item Grouping: extract prompt from conversation context (first 200 characters of chosen response text as proxy)
\end{itemize}

\textbf{Why Anthropic-HH:}
\begin{enumerate}
\item \textbf{Public availability:} Open-source dataset, reproducible validation
\item \textbf{Scale:} 160K+ comparisons provide large pool for prompt sampling
\item \textbf{RLHF relevance:} Used in Constitutional AI, representative of standard evaluation datasets
\item \textbf{Raw data access:} Provides individual comparison records (not just aggregated metrics)
\end{enumerate}

\subsection{Dataset Structure Documentation}

We document two dataset format categories relevant for entropy measurement:

\begin{table}[h]
\centering
\caption{Dataset format categories for entropy measurement.}
\label{tab:dataset-formats}
\begin{tabular}{lp{2.5cm}ccp{2cm}}
\toprule
\textbf{Format Type} & \textbf{Structure} & \textbf{Cost} & \textbf{Entropy} & \textbf{Example} \\
\midrule
Pairwise Unique & Each example: different candidates & Low & No & Anthropic-HH \\
Multi-Annotator Fixed & N annotators: same candidates & High & Yes & OpenAI Sum. \\
\bottomrule
\end{tabular}
\end{table}

\textbf{Key Insight:} Pairwise-unique format optimizes annotation cost by assigning different response pairs to each comparison. This prevents per-prompt entropy aggregation because:
\begin{itemize}
\item \textbf{Intended aggregation:} ``For prompt P, what \% prefer response A vs B vs C?'' (requires fixed \{A, B, C\})
\item \textbf{Actual structure:} ``Comparison 1: chosen=A1, rejected=B1; Comparison 2: chosen=A2, rejected=B2'' (no shared response set)
\item \textbf{Tautological result:} Treating each comparison as ``1 vote for chosen category'' yields 50/50 split (every example has 1 chosen + 1 rejected) $\rightarrow$ H = ln(2) for all prompts
\end{itemize}

\textbf{Multi-annotator-fixed format} avoids this by fixing response candidates per prompt:
\begin{itemize}
\item Example: 5 annotators rate \{Summary\_A, Summary\_B, Summary\_C, Summary\_D\} for article X
\item Votes: \{A: 2, B: 1, C: 1, D: 1\} $\rightarrow$ distribution [0.4, 0.2, 0.2, 0.2] $\rightarrow$ H = 1.33 nats
\item Different article Y has different preference distribution $\rightarrow$ variance across prompts
\end{itemize}

\subsection{Implementation}

\textbf{Tools:}
\begin{itemize}
\item \textbf{Dataset loading:} Hugging Face \texttt{datasets} library (load\_dataset(``Anthropic/hh-rlhf''))
\item \textbf{Entropy computation:} \texttt{scipy.stats.entropy} (handles normalization automatically)
\item \textbf{Visualization:} matplotlib for figures (histogram, scatter plot)
\end{itemize}

\textbf{Validation Checks:}
\begin{itemize}
\item All entropy values in range [0, ln(2)] $\approx$ [0, 0.693] nats
\item No NaN or Inf values
\item Deterministic results (seed=1 reproducibility)
\end{itemize}

\subsection{Expected Outcomes}

\textbf{If H-E1 Passes (A5 valid):}
\begin{itemize}
\item Success rate $\geq$95\% (entropy computable for most prompts)
\item Variance > 0 (entropy values show spread across prompts)
\item Mean entropy 0.3-0.5 nats (moderate diversity, neither uniform nor consensus)
\item Proceed to full hypothesis mechanism tests (H-M1-M4: checkpoint analysis, inflection point detection)
\end{itemize}

\textbf{If H-E1 Fails (A5 violated):}
\begin{itemize}
\item Success rate <95\% OR variance = 0 (constant entropy)
\item Root cause analysis required:
  \begin{itemize}
  \item \textbf{Scenario 1:} Dataset lacks raw counts (only aggregated stats) $\rightarrow$ need new data collection
  \item \textbf{Scenario 2:} Dataset format incompatible (pairwise unique responses) $\rightarrow$ use alternative proxy
  \item \textbf{Scenario 3:} Implementation bug $\rightarrow$ fix and revalidate
  \end{itemize}
\end{itemize}

\textbf{Implications:} H-E1 outcome determines whether hypothesis is testable with existing datasets (pass) or requires methodological pivot (fail).
